% 13_body.tex — Day3 산출물 3/5 (⑬) 이슈 로그 · 리스크 등록부 갱신 (빈 템플릿 / 완성 예시 공용 본문)
% 드라이버: 13_이슈리스크갱신_빈템플릿.tex (\wbblanktrue) / 13_이슈리스크갱신_완성예시.tex (\wbblankfalse)
% 내용 출처: PM트랙/Day3_워크북.md §3.3(항목 가이드) §3.4(빈 템플릿) §3.5(온담 P1 완성 예시본)
\documentclass[10pt,a4paper]{article}
\usepackage[a4paper,margin=16mm,top=14mm,bottom=20mm]{geometry}
\def\DocNum{3}
\def\DocTitle{이슈 로그 · 리스크 등록부 갱신}
\def\DocSub{⑬ Issue Log \& Risk Register Update}
\def\DocVariant{\B{빈 템플릿}{온담 P1 완성 예시}}
\usepackage{wbstyle}
\begin{document}

\docheader
 {\Bg{[정식 명칭과 코드]}{차세대 주문·재고 통합 플랫폼 구축 (ONE ONDAM P1)}}
 {\Bg{[이슈 로그 vN / 등록부 v전→후 / 데이터 일자]}{이슈 로그 v1.0 / 등록부 v1.1 → v1.2 / 데이터 2026-09-30, 발행 10-23}}
 {\Bg{[P1 PM · PMO]}{P1 PM, P1 PMO (소유자 각자 서명)\rd{7mm}}}
 {\Bg{[스폰서 실명]}{정미란 CFO (RACI 26행 A = P1 PM, 게이트 갱신 확인)\rd{7mm}}}

% ------------------------------------------------------------------ 1
\secbar{1. 문서 정보와 갱신 범위}
\guide{이슈 로그 버전·기간, 등록부 버전(전·후)과 갱신 시점(월 갱신 / 게이트), 입력 문서, 이슈 심각도 척도. 심각도는 이슈용(1 = 임계경로·규제·계약·허용범위 초과 영향 / 2 = 허용범위 내 영향·기한 있는 조치 / 3 = 관리 사항)으로 리스크의 P·영향과 구분한다 — 리스크 매트릭스 점수를 이슈 심각도로 쓰지 않는다.}
{\footnotesize
\begin{tabularx}{\linewidth}{|L{30mm}|Y|}\hline
\hd{항목} & \hd{내용} \\ \hline
\B{\rd{9mm}}{}이슈 로그 & \Bg{[vN, 기간, n건 (종결 n · 진행 n) — 결함·변경요청은 별도 문서]}{v1.0, 2026-01-05 ~ 2026-09-30, 11건 (종결 5 · 진행 6). 결함은 ⑭ 결함 관리 도구, 변경요청은 ⑪ 변경 로그에 별도} \\ \hline
\B{\rd{9mm}}{}리스크 등록부 & \Bg{[v전(시점·내용) → v후(데이터 일자, 발행일, 서명 건수)]}{v1.1(G1 2026-04-30 — R-01 분리 의결·R-14 예비 정량화 반영) → \textbf{v1.2}(데이터 2026-09-30, CCB 10-15 결과 반영해 10-23 발행, 소유자 서명 22건)} \\ \hline
\B{\rd{11mm}}{}입력 & \Bg{[성과보고 n호, 변경 로그, 감리 보고, 진단 보고, 변경협의체 합의서, 외부 통보]}{성과보고 제9호(⑫), 변경 로그 CR-01~11(⑪), 감리 1차 보고(09-25), 보안진단 1회차(09-08), 변경협의체 합의서(07-15·09-30), 가맹 부속서 동의 집계(08-14), P2 T0 통보(04-24), 3PL 회신(07-08)} \\ \hline
\B{\rd{11mm}}{}이슈 심각도 & \Bg{[1 = … / 2 = … / 3 = …]}{\textbf{1} = 임계경로·규제 기한·계약(검수·변경협의체)·허용범위 초과에 직접 영향 / \textbf{2} = 허용범위 내 영향, 기한 있는 조치 필요 / \textbf{3} = 관리 사항(기록·회신)} \\ \hline
\B{\rd{9mm}}{}작성 / 승인 & \Bg{[PM, PMO / 스폰서 — RACI 행]}{P1 PM, P1 PMO / 정미란 (RACI 26행 A = P1 PM, 게이트 갱신은 스폰서 확인)} \\ \hline
\end{tabularx}}

% ------------------------------------------------------------------ 2 (가로)
\landscapeon
\secbar{2. 이슈 로그}
\guide{ID(IS-nn), 발생일(인지일 병기), 제목·요약(사실 한 줄 + 영향 한 줄), 출처(리스크 전환 R-nn / 신규 발견 / 감리 / 외부 / 수행사 통보), 심각도, 소유자(실명, 괄호에 조치 실행자), 조치, 기한, 에스컬레이션 단계(1·2·3·병행), 상태(진행 / 종결 + 일자), 연결. \textbf{소유자가 전부 PM이면 틀렸다} — PM이 소유자인 이슈는 PM 권한 안에서 조치 가능한 것뿐. 기한 “ASAP” 금지. 종결은 확인자·일자.}
{\scriptsize\setlength{\tabcolsep}{2pt}
\begin{longtable}{|L{9mm}|L{12mm}|L{50mm}|L{16mm}|C{5mm}|L{22mm}|L{63mm}|L{11mm}|L{18mm}|L{18mm}|L{27mm}|}\hline
\hd{ID} & \hd{발생일 (인지일)} & \hd{제목 · 요약 (사실 / 영향)} & \hd{출처} & \hd{심} & \hd{소유자 (조치자)} & \hd{조치} & \hd{기한} & \hd{에스컬} & \hd{상태} & \hd{연결} \\ \hline \endhead
\ifwbblank
\rd{9mm}IS-01 & & & \g{[R-nn 전환/신규/감리/외부/수행사]} & & \g{[실명 (실명)]} & \g{[무엇을 했고 할 것인가]} & & \g{[1/2/3/병행]} & \g{[진행/종결 일자]} & \g{[R-·CR-·EX-·⑭·⑮]} \\ \hline
\rd{9mm}IS-02 & & & & & & & & & & \\ \hline
\rd{9mm}IS-03 & & & & & & & & & & \\ \hline
\rd{9mm}IS-04 & & & & & & & & & & \\ \hline
\rd{9mm}IS-05 & & & & & & & & & & \\ \hline
\rd{9mm}IS-06 & & & & & & & & & & \\ \hline
\rd{9mm}IS-07 & & & & & & & & & & \\ \hline
\rd{9mm}IS-08 & & & & & & & & & & \\ \hline
\rd{9mm}IS-09 & & & & & & & & & & \\ \hline
\rd{9mm}IS-10 & & & & & & & & & & \\ \hline
\rd{9mm}IS-11 & & & & & & & & & & \\ \hline
\rd{9mm}IS-12 & & & & & & & & & & \\ \hline
\else
IS-01 & 02-20 & \textbf{R-04 실현} — CFO 1차연도 148억 확정, P1 2026년 배분 102.0 → 94.0(−8.0). 영향: 4월 라이선스 2차(6.0)·5월 SAP 애드온(3.4) 지급 불가 & R-04 전환 & 1 & 정미란 (재경팀장, 강현도) & 라이선스 2·3차분 지급 조건 “인도 50\% / 통합시험 통과 50\%” 계약 변경(03-13), 인도 5·7월 재조정, A05 개발환경은 선약정으로 영향 없음 확인. 관리예비비 요청 없이 흡수. PV 곡선 유지(재기준선 아님) & 03-13 & 2단계(스폰서 결정 03-06, 10영업일 내) & \textbf{종결} 03-13 (재경팀장 확인) & R-17 등재, ⑫ 착시 1, ⑮ 계약 변경 1번 \\ \hline
IS-02 & 05-12 (05-14) & \textbf{R-05 실현} — POS 화면 정의서 v1.0 승인 대기 12영업일(오정근 부속서 협상 집중, 대체 결정권자 미지정). S10 3주로 연장, 이후 스프린트 1주 순연 & R-05 전환 & 1 & 박세연 (오정근, 매장운영팀장) & 대체 결정권자 지정(매장운영팀장, 06-03 — 3주 늦음), 결정 SLA 5영업일 재공지, 변경협의체 07-15 판정: 대기 12영업일 중 발주사 귀책 7일, 대기 비용 0.42억(컨틴전시 R-05 집행) & 07-15 & 병행(변경협의체 05-27 상정) & \textbf{종결} 07-15 — 재발 리스크 R-05 재평가 유지 & CR-06, ⑫ 신호 2, ⑮ 대기 비용 \\ \hline
IS-03 & 03-18 & \textbf{인터페이스 13본 실체 확인} — 문서 부재 OD-POS↔SAP 6본 중 2본은 2019년 이후 미사용(폐기), 목록에 없던 2본 발견(OD-POS↔멤버십 정산 일배치, 가맹발주웹↔BI 추출). 총 47본 유지(45 + 2), 설계문서 대상 47본 & 신규 발견 (A02 복원 중, 박세연) & 2 & 박세연 (수행 인터페이스 PL) & 인터페이스 목록 v2 확정(04-03), 신규 2본 설계문서·재구축 FP +40 → CR-07(09-17 PM 전결), 폐기 2본은 컷오버 시 차단 목록 & 09-17 & 1단계 & \textbf{종결} 09-17 & CR-07, HR-07 문서 47본 \\ \hline
IS-04 & 08-31 & \textbf{슈퍼유저 이탈} — P3 슈퍼유저 124명 중 11명 이탈(퇴사 6·전보 5, 강윤지 지역매니저 포함). 영향: 설계 검토회 매장 대표 공백, UAT 참여자 재선정, 파일럿 12개점 중 2개점 교체 필요 & 외부(P3) & 2 & P3 PM (오정근) & 재선정 11-13 통보, P1은 UAT TC 리뷰를 12월로, 파일럿 후보 교체 시 구형 기종 2점 포함 조건 유지, 강윤지 후임 지역매니저 검토회 초청 & 11-13 & 병행(운영위 09-10 보고) & 진행 & CR-08 조건 ④, R-13, R-15 \\ \hline
IS-05 & 06-30 (07-08) & \textbf{R-03 트리거 발동} — 대성로지스 협상 06-30 미타결, 07-08 한동석 회신 “준실시간 불가, 야간 배치 2회 폴백 확정, 2027-06 재협상 조항” & R-03 전환 & 2 & 한동석 (수행 인터페이스팀) & 폴백 모드 확정, 준실시간 API는 명세·구현 완료 후 보관(WBS 1.3.3 카드 외부 의존 조항대로), 폴백 자동 전환 시험 08-21 통과(TC-INV-031~040), 부하 산정서 갱신. 컨틴전시 R-03 0.50 집행 — \textbf{8월 CCB 보고 누락, 10-15 소급 보고} & 08-21 & 병행(소관 임원 → 스폰서 07-10) & \textbf{종결} 08-21 — R-16 정식 등재 & R-16, IS-07(감리 8번), ⑫ A08 TF 20 → 12 \\ \hline
IS-06 & 08-14 & \textbf{가맹 부속서 동의 71\%}(276/388). 잔여 112점 중 비호환 단말 32점 겹침. 영향: 가맹 SW 배포 범위(1.1.6) 미확정 → 배포 패키지 R4 착수 조건 & 외부(P3·오정근) & 2 & 오정근 (매장운영팀장) & 잔여 동의 12월 목표(선행 릴리스 12-04를 지렛대로, O-02), 배포 범위 확정 시한 G2 11-27, 미동의 매장은 컷오버 후 2차 배포 옵션(WBS 1.1.6 대응안 ③), 협의회 22:00 요구(CR-09)를 관계 개선에 활용 & 11-27 & 병행(소관 임원) & 진행 & R-02, R-13, CR-09, ⑭ 감리 4번 \\ \hline
IS-07 & 09-25 & \textbf{감리 1차 지적 8건} 통보. 그중 3번(P2 인터페이스 규격 설비측 확인 부재)·4번(가맹 SW 배포 범위 설계 미확정)이 “수행사 설계 미흡”으로 기재 — 실제 원인은 발주사(프로그램 PMO·오정근) 미결정. 영향: R2 검수 보류, 수행사 대기 비용 청구 근거 & 감리 & 1 & \textbf{P1 PM} (정하늘, 이재현) & 이의 제기 서면 09-29(⑭ 대응 관리대장), 귀속 재분류 회신 요청 10-08, 발주사 귀속 2건은 발주사 조치(P2 규격 확인 10-16, 배포 범위 11-27), 수행사 귀속 3건 조치 10-16~10-30, R2 검수 보류 서면 회신(⑮) & 10-08 & 2단계(스폰서 회신 요청 09-30) & 진행 & ⑭ 항목 6, ⑮ 검수 보류, R-20 \\ \hline
IS-08 & 09-08 & \textbf{보안진단 1회차 High 2건} — API GW 인증 토큰 만료 미검증, MDM 관리 콘솔 접근 통제 미비(Medium 9건 별도). 지속가능성 축 “High 미조치 0(G3)” & 3자 진단 & 1 & 박세연 (수행 통합팀, API GW 벤더) & 조치 10-16, 재진단 10-23, 최영주 보고(09-10), 감리 6번(파기 주기 증적)과 함께 규제 매핑표 REG-01 갱신 & 10-16 & 병행(최영주) & 진행 & 격자 지속가능성, ⑭ 감리 6번 \\ \hline
IS-09 & 09-17 & \textbf{서정필 22:00 마감 요구} 서면 접수(협의회 공문 2026-14호) & 외부(협의회) & 3 & 오정근 (P1 PM) & CR-09 영향 분석·변경협의체 09-30·CCB 10-15 조건부 승인(파라미터), 협의회 회신 10-23 & 10-23 & — & 진행 (CR-09) & CR-09, R-16, IS-06 \\ \hline
IS-10 & 07-31 & \textbf{클라우드 개발환경 사용량 초과} — 월 0.4 → 0.5(시험 환경 2벌 상시 가동, 6월부터). 영향: 인프라 CV −0.60 누적, 유보 3.0 복귀 가능성 & 신규 발견 (AC 대조, 재경) & 3 & 박세연 (강현도, 수행 통합팀) & 야간·주말 자동 축소 정책 09-01 적용(효과 10월 청구서로 확인), 사이징 재검토는 2027-01 부하시험 후, 유보 3.0 복귀 판단 시점 명시 & 10-31 & 1단계 & 진행 & ⑫ 인프라 CV, Day2 ⑨ 조정표 “인프라 유보” \\ \hline
IS-11 & 07-20 & \textbf{수행사 인터페이스 PL 교체} 통보(개인 사유 퇴직). 1.3.3(3PL)·1.5.5(47본) 핵심 인력 & 수행사 통보 & 2 & 이재현 (P1 PM 승인) & 후임 이력·경력 검토(07-22), 2주 인수인계(07-20~08-02), 승인 조건: 3PL 명세·13본 복원 지식 이전 체크리스트 확인(박세연 08-03 서명), 교체 승인서(⑮ 항목 3) & 08-02 & 1단계 & \textbf{종결} 08-03 & ⑮ 수행사 성과 기록, R-06 유형 \\ \hline
\fi
\end{longtable}}
\B{\small\g{출처 분포: 리스크 전환 [ ] / 신규 발견 [ ] / 감리·진단 [ ] / 외부 [ ] / 수행사 [ ] — 등록부에 없던 원인 중 프리모템 원문에 있었던 것: \hrulefill}\par}{\small 출처 분포: 리스크 전환 3(IS-01·02·05) / 신규 발견 2(IS-03·10) / 감리·진단 2(IS-07·08) / 외부 3(IS-04·06·09) / 수행사 1(IS-11). \textbf{등록부에 없던 원인 5건(IS-03·04·08·10·11) 중 3건(IS-04·08·11)은 프리모템 원문(6번 개발자 이탈, 12번 보안, 13번 교육 이탈)에 있었으나 리스크로 승격되지 않았던 것} — Day4 교훈 등록부(⑲) 입력.\par}

% ------------------------------------------------------------------ 3 (가로)
\secbar[55mm]{3. 리스크 → 이슈 전환 규칙}
\guide{“트리거 발동”과 “실현”을 구분한다 — 트리거는 신호이고 실현은 사건이다. 실현(P 100\%)된 리스크를 등록부에 그대로 두지 않는다. 트리거 발동을 실현으로 오해해 컨틴전시를 미리 쓰지 않는다. 부분 실현의 처리와 관리예비비 대상 리스크의 즉시 보고를 규칙으로 적는다.}
\begin{fieldbox}\small
\ifwbblank
\g{트리거 발동 = 신호 → [기록·집행 준비] / 실현 = P 100\% → [IS-nn 등재, 컨틴전시 배정 집행 개시] / 잔여 있으면 [같은 번호로 재평가] / 부분 실현 → [ ] / 대응 완료·종결 → [새 노출 등재] / 관리예비비 대상 트리거 → [당일 보고]}\lines{5}{7mm}
\else
\textbf{리스크 관리 계획 v1.0 §4 [추가 설정]}\par
• \textbf{트리거 발동}은 신호다 — 소유자가 관측 사실과 일자를 등록부에 기록하고, 조건부 대응(컨틴전시 배정분)의 집행 준비를 시작한다. 발동 자체로 컨틴전시를 쓰지 않는다.\par
• \textbf{실현}은 사건이다 — 확률 100\%가 된 순간 등록부 행을 “실현 → IS-nn”으로 바꾸고 이슈 로그에 등재하며, 컨틴전시 배정분 집행을 개시한다(RACI 19행). 실현 후에도 같은 원인의 재발 가능성이 있으면 같은 번호로 잔여 리스크를 재평가해 유지한다(R-05).\par
• \textbf{부분 실현}(R-04: 삭감 확정, 영향은 대응으로 흡수)은 실현으로 기록하고, 흡수 조치가 만든 새 노출(R-17)을 2차 리스크로 등재한다.\par
• \textbf{대응 완료·종결}(R-03: 폴백 확정) — 사건은 확정되었으나 원가 영향이 배정분 안에서 끝났으면 종결하고, 폴백이 만드는 새 노출(R-16)을 등재한다.\par
• \textbf{관리예비비 대상 리스크(R-01·R-04·R-14·R-17)의 트리거 발동은 당일 스폰서·프로그램 PMO장 보고}(Day2 ⑩ 항목 9). 컨틴전시 집행은 월 CCB 보고 — \textbf{누락 시 감리 지적 대상(IS-05·감리 8번)}.
\fi
\end{fieldbox}

% ------------------------------------------------------------------ 4 (가로)
\clearpage
\secbar{4. 리스크 재평가 (v\B{[전]}{1.1} → v\B{[후]}{1.2})}
\guide{등록부의 \textbf{모든 행}에 대해 — 상태(유지 / 트리거 발동 / 실현 → IS / 종결 / 이관), 갱신 전 잔여(P × 영향 = EMV) → 갱신 후, 변경 근거(관측된 사실·일자), 자금원 변동, 다음 트리거·검토 시점. 값이 바뀌지 않은 행도 “유지 — 근거”를 적는다(검토했다는 증거). \textbf{확률을 낮추는 갱신에는 관측 사실을 반드시 붙인다.} 종결에는 종결 사유. 나쁜 리스크만 갱신하지 않는다.}
{\scriptsize\setlength{\tabcolsep}{2pt}
\begin{longtable}{|L{20mm}|L{30mm}|L{30mm}|L{24mm}|L{88mm}|L{28mm}|L{32mm}|}\hline
\hd{ID} & \hd{상태} & \hd{잔여 전 (P×영향=EMV)} & \hd{잔여 후} & \hd{근거 (사실 · 일자)} & \hd{자금원} & \hd{다음 트리거 · 검토} \\ \hline \endhead
\ifwbblank
\rd{7mm}R-01 & \g{[유지/발동/실현→IS/종결/이관]} & \g{[P × 영향 = EMV]} & & \g{[관측 사실, 일자]} & \g{[컨틴전시/관리예비비/해제]} & \\ \hline
\rd{7mm}R-02 & & & & & & \\ \hline
\rd{7mm}R-03 & & & & & & \\ \hline
\rd{7mm}R-04 & & & & & & \\ \hline
\rd{7mm}R-05 & & & & & & \\ \hline
\rd{7mm}R-06 & & & & & & \\ \hline
\rd{7mm}R-07 & & & & & & \\ \hline
\rd{7mm}R-08 & & & & & & \\ \hline
\rd{7mm}R-09 & & & & & & \\ \hline
\rd{7mm}R-10 & & & & & & \\ \hline
\rd{7mm}R-11 & & & & & & \\ \hline
\rd{7mm}R-12 & & & & & & \\ \hline
\rd{7mm}R-13 & & & & & & \\ \hline
\rd{7mm}O-01 & & \g{(−)} & & & & \\ \hline
\rd{7mm}O-02 & & \g{(−)} & & & & \\ \hline
\else
R-01 P2 리드타임 & \textbf{트리거 발동 → 대응 실행 → 잔여 재평가} & 60 × 2.0 = 1.20 (v1.1: 분리 의결 후 40 × 2.0 = 0.80) & \textbf{30 × 2.0 = 0.60} & T0 04-24 리드타임 38주 통보 → MC 2027-03-05. G1 이사회 의결 “매장·주문 1차 02-26 / 물류 2차 03-27”. 잔여 = MC 추가 지연 시 2차 컷오버 창 이탈·재시험. P2 9월 진도 보고 정상 & 관리예비비 & 2026-10 P2 진도, 2027-01 MC 확정 통보 \\ \hline
R-02 가맹 동의 & 유지 (트리거 미발동) & 30 × 2.0 = 0.60 & \textbf{40 × 1.5 = 0.60} & 7월 중간 집계 64\%(트리거 60\% 미만 아님), 08-14 최종 71\%. 잔여 112점 중 비호환 32점 겹침 → 확률 상향, 영향은 분리 배포 준비(1.1.6 대응안)로 하향 & 컨틴전시 0.60 & 12월 동의율, G2 배포 범위 확정 \\ \hline
R-03 3PL 거부 & \textbf{트리거 발동 → 대응 완료 → 종결} & 40 × 1.2 = 0.48 & \textbf{종결 (0)} & 06-30 미타결, 07-08 폴백 확정 회신, 08-21 폴백 시험 통과. 배정 0.50 전액 집행. 새 노출은 R-16 & 집행 0.50 & — (2027-06 재협상은 P1 밖) \\ \hline
R-04 1차연도 삭감 & \textbf{실현 → IS-01} & 40 × 2.5 = 1.00 & \textbf{실현 (0)} & 02-20 확정, 지급 조건 변경으로 흡수, 관리예비비 미사용. 새 노출 R-17 & — & — \\ \hline
R-05 겸임 인력 대기 & \textbf{실현 → IS-02 → 재발 리스크 유지} & 40 × 1.5 = 0.60 & \textbf{50 × 1.0 = 0.50} & 05-12 실현, 대기 비용 0.42 집행. 대체 결정권자 지정(06-03) 후 영향 하향, G2 정본 판정·UAT 결정 집중으로 재발 확률 상향 & 컨틴전시 잔여 0.18 + 재배정 0.32 = 0.50 & 월 참석률, 10-19 정본 판정 착수 \\ \hline
R-06 OD-POS 개발자 & 하락 & 20 × 1.5 = 0.30 & \textbf{10 × 1.5 = 0.15} & 13본 복원 03-25 완료(O-01 실현), 지식 이전 세션 기록 6회, 인사본부 유지 협의 완료. 잔여 = 전환 리허설·CR-08 임시 연동 시 필요 & 컨틴전시 0.20 (0.10 해제) & 2026-12 A12 착수 \\ \hline
R-07 동결 직전 집중 & \textbf{트리거 발동 → 대응 실행 → 유지} & 40 × 1.0 = 0.40 & \textbf{40 × 1.0 = 0.40} & FZ-05 순변동 +41 > 35 발동, FP 재산정(CR-06) 0.15 집행, 22건 기각. FZ-06~09는 마감 규칙 정착으로 순변동 ±20 이내 & 컨틴전시 잔여 0.25 + 0.15 = 0.40 & 매월 동결 \\ \hline
R-08 앱 벤더 & 하락 & 30 × 2.0 = 0.60 & \textbf{25 × 2.0 = 0.50} & 소스 인수·API 연동 조항 계약 변경 04-22 서명, API 규격 04-30 제공, 벤더 연동 착수 07. 잔여 = 12월 연동 시험 실패. 단가 인상 요구는 R-18 & 컨틴전시 0.50 (0.10 해제) & 12월 연동 시험 \\ \hline
R-09 별도 법인 규제 & \textbf{종결} & 20 × 1.5 = 0.30 & \textbf{종결 (0)} & P4 위탁 구조 05-20 확정, 온담푸드시스템 위탁 계약 06-12 서명(최영주·조성태), 공장 연동 1.3.5 설계 반영 & 0.30 해제 & — \\ \hline
R-10 G2 후 마스터 변경 & \textbf{상승} & 10 × 3.8 = 0.38 & \textbf{20 × 3.8 = 0.76} & 신규 브랜드(델리 서브 브랜드) 상품 코드 요구 1건·B2B 거래처 코드 체계 변경 요구 1건 접수(09월, 영업본부·구매팀) — 정본 판정 전 반영 협의 중. 판정 후 변경 요구 재발 확률 상향 & 컨틴전시 0.40 → \textbf{0.80} (미배정 풀에서 0.40) & 10월 접수 마감, G2 동결 규칙 의결 \\ \hline
R-11 SAC 미합의 & 하락 & 20 × 2.0 = 0.40 & \textbf{15 × 2.0 = 0.30} & SAC 초안 G1 합의(RACI 22행, 04-24), 운영 요건 월 검토 5회. 잔여 = G3 사전 점검 미해소 항목 & 컨틴전시 0.30 (0.10 해제) & G3 사전 점검 2027-01 \\ \hline
R-12 편익 분모 & \textbf{종결·이관} & 25 × 0.8 = 0.20 & \textbf{이관 (0)} & M+3(04-28) 두 분모 병기 합의(나영선·한동석), 리포트 StRS-034 R3 인도 예정. 실사 협조는 편익 실현 원장(Day4 ⑳)에서 관리 & 0.20 해제 & G5 \\ \hline
R-13 구형 단말 & 유지 & 15 × 2.1 = 0.32 & \textbf{15 × 2.1 = 0.32} & 04-30 조사 비호환 32점(8.2\%) — 트리거 10\% 미발동. 파일럿에 구형 2점 포함 유지 & 컨틴전시 0.30 & 2027-01 파일럿 배포 결과 \\ \hline
O-01 문서화 조기 완료 & \textbf{실현 → 종결} & 50 × −0.8 = −0.40 & \textbf{종결 (0)} & 03-25 완료(80\% 이상 03-06 관측), A02 정시, 전환 매핑 +0.20 선행 & — & — \\ \hline
O-02 포털 선행 릴리스 & 상승(기회) & 60 × −1.0 = −0.60 & \textbf{70 × −1.0 = −0.70} & 05-20 프로토타입 시연 협의회 반응 긍정(설문 4.2/5), 선행 릴리스 12-04 예정. CR-09 파라미터 기능이 지렛대 추가 & — & 12-04 선행 릴리스 \\ \hline
\fi
\end{longtable}}
\B{}{\small 15행 갱신: 상승 2(R-10, R-02 영향 재구성) · 하락 5(R-01·R-05·R-06·R-08·R-11) · 종결 4(R-03·R-09·R-12·O-01) · 실현 2(R-04·R-05) · 유지 2(R-07·R-13).\par}

% ------------------------------------------------------------------ 5 (가로)
\clearpage
\secbar{5. 2차 · 신규 리스크 정식 등재}
\guide{Day2에서 예비 번호로 둔 2차 리스크(R-14~R-18)와 실행 중 새로 식별된 리스크(R-19·R-20)를 등록부 본표와 같은 열로 정량화한다. 출처(어느 대응·이슈·성과 분석에서 나왔는가)를 적는다. 관리예비비 대상 여부는 \textbf{원인의 계층}으로 판정한다. 새 리스크의 소유자를 PM으로 두는 것이 흔한 오류다.}
{\scriptsize\setlength{\tabcolsep}{2pt}
\begin{longtable}{|L{9mm}|L{50mm}|L{18mm}|C{8mm}|L{18mm}|C{10mm}|L{60mm}|L{22mm}|L{38mm}|L{20mm}|}\hline
\hd{ID} & \hd{원인 → 사건 → 영향} & \hd{출처} & \hd{P\%} & \hd{영향 (억 / 영업일)} & \hd{EMV} & \hd{전략 · 조치 (WBS, 기준선 포함 여부)} & \hd{소유자 (조치자)} & \hd{트리거} & \hd{자금원} \\ \hline \endhead
\ifwbblank
\rd{11mm}R-14 & \g{[원인 → 사건 → 영향]} & \g{[R-nn 대응 / IS-nn / ⑫ / ⑭]} & & \g{[억 / 일]} & \g{[P×억]} & \g{[전략 — 조치, WBS 코드, 기준선 포함 여부]} & \g{[실명 (실명)]} & \g{[관측 가능 신호·시점]} & \g{[컨틴전시 n / 관리예비비]} \\ \hline
\rd{11mm}R-15 & & & & & & & & & \\ \hline
\rd{11mm}R-16 & & & & & & & & & \\ \hline
\rd{11mm}R-17 & & & & & & & & & \\ \hline
\rd{11mm}R-18 & & & & & & & & & \\ \hline
\rd{11mm}R-19 & & & & & & & & & \\ \hline
\rd{11mm}R-20 & & & & & & & & & \\ \hline
\else
\textbf{R-14} & R-01 분리 대응 → 매장·주문(02-26)·물류(03-27) 두 번 컷오버 → 매장 혼선·재리허설·2차 창 성수기 접근(설 이후 3월 B2B 납품 집중 [추가 설정]) & R-01 대응 (G1 의결) & 50 & 1.6 / 10 (2차 창) & \textbf{0.80} & 완화 — 컷오버 런북(Day4 ⑯)을 2회 분리 작성, 1차 하이퍼케어 지표(심각도 2 이상 5건 이하)를 2차 go 조건으로, 매장 통보 1회로 통합(1.6.9 기준선 포함, 추가 FP 0) & 프로그램 PMO장 (P1 PM, 박준영) & 1차 컷오버 후 하이퍼케어 2주 심각도 2 이상 5건 초과; P2 MC 2027-03-05 지연 통보 & \textbf{관리예비비} (원인: 프로그램 리드타임·버퍼) \\ \hline
\textbf{R-15} & 포털 선행 릴리스(12-04) 결함 → 협의회 신뢰 하락 → 잔여 112점 동의 협상 지렛대 상실 & R-02·O-02 공유 & 30 & 0.8 / 0 & \textbf{0.24} & 완화 — 선행 릴리스 범위를 발주·마감·배송 추적 3기능으로 한정(1.4.3), 협의회 지정 점주 5명 사전 검수(범위 기술서 §4 ④-③), CR-09 파라미터 포함 & 오정근 (수행 포털팀) & 사전 검수 심각도 2 이상 3건 초과(11-27) & 컨틴전시 0.30 \\ \hline
\textbf{R-16} & R-03 폴백(야간 배치 21:00·03:00) + 22:00 마감(CR-09) → 야간 작업 창 충돌 → 배치 지연·재고 정합 오류 & R-03 대응, CR-09 & 30 (CR-09 22:00 운용 시 45) & 0.8 / 5 & \textbf{0.24} (→ 0.36) & 완화 — 배치 1회차를 22:30으로 조정 가능하게 스케줄 파라미터화(1.3.3, FP 0 — 기존 설계 범위), 부하 산정서 v2에 22:00 시나리오 추가 & 한동석 (수행 인터페이스팀) & 폴백 배치 1회차 실행 시간 40분 초과 2회 연속(통합시험) & 컨틴전시 0.30 \\ \hline
\textbf{R-17} & R-04 대응(라이선스 4.70 지급 이연 + 2026 배분 −8.0의 2027 이월) → 2027 자금 계획 초과 → 2027-01 단말·DR·시험 봉우리(16.20) 집행 지연 & IS-01 & 40 & 2.0 / 10 (임계, A16·단말 설치) & \textbf{0.80} & 완화 — 2027 자금 계획(재무본부 11월 확정)에 이월분 12.7(4.70 + 8.0)을 명시 반영 요청, 단말 발주 10-30 선집행으로 봉우리 분산 & 정미란 (재경팀장) & 2027 자금 계획 초안(11월)에 이월분 미반영 & \textbf{관리예비비} (원인: 스폰서 자금 결정) \\ \hline
\textbf{R-18} & R-08 계약 변경(소스 인수·API 조항) → 벤더가 유지보수 단가 인상(월 3,200 → 4,000만 원 [추가 설정]) 요구 → 협상 교착 시 앱 연동 시험 지연 & R-08 대응 & 30 & 0.5 / 5 & \textbf{0.15} & 완화 — 단가 협상은 앱 유지보수 계약(오정근·박세연) 소관으로 분리, P1은 API 연동 시험 일정(12월)만 계약 변경 조항으로 고정 & 박세연 (오정근) & 벤더 서면 인상 요구(10월 예상) & 컨틴전시 0.20 (P1측 재작업만) \\ \hline
\textbf{R-19} & R2 이월 420 FP + CR-08 90 → R4 잔여(A12) 30일·TF 0 + A14·A16 압축 → 통합시험 범위 축소 → 결함 UAT·운영 유출 & ⑫ 회복 계획 EX-01 & 40 & 1.5 / 10 (임계) & \textbf{0.60} & 완화 — 회귀 자동화 도구·시험 인력 선투입(0.60, 컨틴전시), 통합시험 TC 400건 중 Must 경로 우선 실행 순서 확정(1.6.4), 심각도 1·2 미해결 0 조건 유지 & P1 PM (수행 시험팀, 이재현) & S22 종료(11-13) 속도 620 미만 → 대안 3 발동; 통합시험 2주차 결함 발견율 계획 대비 150\% 초과 & 컨틴전시 0.60 \\ \hline
\textbf{R-20} & 감리 지적 귀속 분쟁(IS-07) → R2 검수 보류 장기화 → 수행사 대기 비용 청구·변경협의체 판정 → 검수 지연이 G2 감리 커버리지에 재지적 & ⑭ 감리 대응 & 30 & 1.0 / 5 & \textbf{0.30} & 완화 — 검수 보류 서면 회신(⑮)에 발주사 사유 기간 명시, 귀속 재분류 회신 10-08 시한, 검수 재개 조건(수행사 귀속 3건 조치) 명시 & P1 PM (정하늘, 이재현) & 10-08 회신 미도착; 변경협의체 10-14 대기 비용 상정 & 컨틴전시 0.30 \\ \hline
\fi
\end{longtable}}

% ------------------------------------------------------------------ 6 (가로)
\secbar[70mm]{6. 잔여 EMV 재계산과 허용범위 대조}
\guide{갱신 후 위협 합·기회 합·순, 자금원별(컨틴전시 대상 / 관리예비비 대상), 단일 최대, 허용범위(8억·단일 2억)와의 대조, 전 버전 대비 변동의 \textbf{분해}(종결 −, 신규 +, 재평가 ±). 일정 EMV(임계경로 직접 영향)도 다시 계산해 잔여 버퍼와 대조한다 — 합계만 적고 분해가 없으면 무엇이 움직였는지 모른다.}
{\footnotesize\setlength{\tabcolsep}{3pt}
\begin{tabularx}{\linewidth}{|L{30mm}|C{18mm}|C{16mm}|C{14mm}|L{28mm}|L{38mm}|L{22mm}|C{16mm}|Y|}\hline
\hd{구분} & \hd{위협 합} & \hd{기회 합} & \hd{순} & \hd{컨틴전시 대상} & \hd{관리예비비 대상} & \hd{단일 최대} & \hd{허용범위} & \hd{판정} \\ \hline
\ifwbblank
\rd{8mm}v\g{[전]} (\g{[일자]}) & & & & \g{[n건]} & \g{[R-ID]} & & 8 / 2 & \\ \hline
\rd{8mm}\textbf{v\g{[후]}} (\g{[일자]}) & & & & & & & 8 / 2 & \g{[이내/초과 — 여유, 예외 보고 여부]} \\ \hline
\rd{11mm}변동 분해 & \multicolumn{8}{l|}{\g{전 [ ] − 종결·실현·이관 [ ] + 신규 [ ] ± 재평가 [ ] = 후 [ ]}} \\ \hline
\rd{8mm}\g{[조건부 시나리오]} & & & & & & & & \\ \hline
\else
v1.0 (2026-03-27) & 6.78 & −1.00 & 5.78 & 4.58 (11건) & 2.20 (R-01·R-04) & R-01 1.20 & 8 / 2 & 이내 (여유 1.22) \\ \hline
\textbf{v1.2 (2026-09-30)} & \textbf{7.26} (16건) & \textbf{−0.70} (O-02) & \textbf{6.56} & \textbf{5.06} (13건) & \textbf{2.20} (R-01 0.60·R-14 0.80·R-17 0.80) & R-14·R-17 0.80 & 8 / 2 & \textbf{이내 (여유 0.74) — 예외 보고 불요, 경고} \\ \hline
변동 분해 & \multicolumn{8}{>{\raggedright\arraybackslash}p{\dimexpr\linewidth-30mm-4\tabcolsep-3\arrayrulewidth\relax}|}{6.78 − 종결·실현·이관 1.98 (R-03 0.48 + R-04 1.00 + R-09 0.30 + R-12 0.20) + 신규 3.13 (R-14 0.80 + R-15 0.24 + R-16 0.24 + R-17 0.80 + R-18 0.15 + R-19 0.60 + R-20 0.30) + 재평가 순증 −0.67 (R-01 −0.60, R-02 0, R-05 −0.10, R-06 −0.15, R-08 −0.10, R-10 \textbf{+0.38}, R-11 −0.10, R-13 0) = \textbf{7.26} (검산 일치)} \\ \hline
CR-09 22:00 운용 시 & 7.31 (R-16 +0.12, R-02 −0.07) & −0.70 & 6.61 & 5.11 & 2.20 & 0.80 & 8 / 2 & 이내 (여유 0.69) \\ \hline
\fi
\end{tabularx}}
\B{\small\g{여유가 줄어든(늘어난) 이유 한 문단 — 종결 대비 신규, 대응이 만든 2차 리스크, 관리예비비 대상의 내용 변화: \hrulefill}\par}{\small \textbf{여유가 1.22에서 0.74로 줄어든 이유}는 종결 4건(−1.98)보다 신규 7건(+3.13)이 크기 때문이고, 신규 7건 중 5건(R-14·15·16·17·18)은 \textbf{대응 조치가 만든 2차 리스크}다. 대응은 노출을 없애지 않고 다른 자리로 옮긴다. 관리예비비 대상 2.20은 v1.0과 같은 숫자이지만 내용이 바뀌었다 — R-04(자금 삭감)가 R-17(자금 이월)로, R-01의 절반이 R-14(두 번 컷오버)로. 셋 다 원인이 프로그램·스폰서 계층이다.\par}
\vspace{1mm}
{\footnotesize
\begin{tabularx}{\linewidth}{|L{30mm}|Y|L{78mm}|L{60mm}|}\hline
\hd{일정} & \hd{임계경로 직접 잔여 영향 (영업일, P × 일수)} & \hd{잔여 버퍼} & \hd{판정} \\ \hline
\ifwbblank
\rd{9mm}v\g{[전]} & \g{[R-nn … = 합]} & \g{허용범위 [ ] + 합류 버퍼 [ ] = [ ]} & \\ \hline
\rd{9mm}\textbf{v\g{[후]}} & & \g{[− 회복 계획 후 여유]} & \\ \hline
\else
v1.0 & 19.0 (R-01·R-04·R-05·R-10·R-11) & 허용범위 15 + 합류 버퍼 15 = 30 & 이내 (여유 11) \\ \hline
\textbf{v1.2} & R-01 3.0 + R-05 2.5 + R-10 4.0 + R-11 1.5 + R-14 5.0 + R-17 4.0 + R-19 4.0 + R-20 1.5 = \textbf{25.5} & 허용범위 15 + 합류 버퍼 15(\textbf{배정 미정}, EX-01 대안 2가 전부 소비) − 회복 계획 후 G3 앞 여유 0 = \textbf{15~30} & \textbf{경고} — 회복 계획이 승인되어도 G3 앞 여유 0이므로 임계경로 리스크 하나 실현 시 합류 버퍼 배정 필요 \\ \hline
\fi
\end{tabularx}}

% ------------------------------------------------------------------ 7 (가로)
\clearpage
\secbar{7. 컨틴전시 소진 현황과 재배정 (컨틴전시 \B{[ ]}{10.20})}
\guide{배정표(Day2 ⑨ 항목 6)의 리스크별 배정·집행(약정과 AC 발생을 구분)·해제·조정·재배정, 미배정 풀의 사용(승인 변경 편입)과 잔여, 관리예비비 요청, 소진율 vs 진척률. \textbf{검산: 재배정 합 + 미배정 풀 + 집행 + BL 편입 = 컨틴전시.} 종결 리스크의 배정이 남아 있으면 틀린 것이다.}
{\footnotesize\setlength{\tabcolsep}{3pt}
\begin{longtable}{|L{24mm}|C{18mm}|L{58mm}|L{62mm}|L{44mm}|L{48mm}|}\hline
\hd{리스크} & \hd{배정 v\B{[전]}{1.0}} & \hd{집행 (약정 / AC 발생)} & \hd{해제 · 조정} & \hd{배정 v\B{[후]}{1.2}} & \hd{근거} \\ \hline \endhead
\ifwbblank
\rd{4mm}R-02 & & & & & \\ \hline
\rd{4mm}R-03 & & & & & \\ \hline
\rd{4mm}R-05 & & & & & \\ \hline
\rd{4mm}R-06 & & & & & \\ \hline
\rd{4mm}R-07 & & & & & \\ \hline
\rd{4mm}R-08 & & & & & \\ \hline
\rd{4mm}R-09 & & & & & \\ \hline
\rd{4mm}R-10 & & & & & \\ \hline
\rd{4mm}R-11 & & & & & \\ \hline
\rd{4mm}R-12 & & & & & \\ \hline
\rd{4mm}R-13 & & & & & \\ \hline
\rd{4mm}\g{[신규 R-nn]} & — & & \g{[+배정]} & & \g{[신규]} \\ \hline
\rd{4mm}\g{[신규 R-nn]} & — & & & & \\ \hline
\rd{4mm}\g{[신규 R-nn]} & — & & & & \\ \hline
\rd{4mm}\textbf{배정 소계} & & & & \g{[n건]} & \\ \hline
\rd{6mm}미배정 풀 & & \g{[CR-nn 발생 / 약정]} & \g{[승인 변경 BL 편입 −, 배정 이동 −]} & \g{[잔여 (스프린트 n회분)]} & \\ \hline
\rd{6mm}\textbf{합계} & & & & \g{[배정 + 미배정 + 집행 + BL 편입 = 컨틴전시]} & \\ \hline
\rd{6mm}관리예비비 요청 & \multicolumn{5}{l|}{\g{[n건, 금액 / 잔여 / 대상 리스크 EMV / 유보 복귀 조건]}} \\ \hline
\rd{6mm}소진율 & \multicolumn{5}{l|}{\g{약정 [ ] ÷ 컨틴전시 [ ] = [ ]\% vs 원가 진척 EV/BAC = [ ]\% — 판독}} \\ \hline
\else
R-02 & 0.60 & — & 0 & 0.60 & 유지 \\ \hline
R-03 & 0.50 & \textbf{0.50 / 0.50} (08월, 폴백 재시험·부하 산정) & −0.50 → 0 & 0 (종결) & IS-05 \\ \hline
R-05 & 0.60 & \textbf{0.42 / 0.42} (07월, 대기 비용) & 잔여 0.18 + 0.32 & 0.50 & 재발 리스크 재평가 \\ \hline
R-06 & 0.30 & — & −0.10 & 0.20 & 하락 \\ \hline
R-07 & 0.40 & \textbf{0.15 / 0.15} (06월, FP 재산정 공수) & 잔여 0.25 + 0.15 & 0.40 & 유지 \\ \hline
R-08 & 0.60 & — & −0.10 & 0.50 & 하락 \\ \hline
R-09 & 0.30 & — & −0.30 & 0 (종결) & 해제 \\ \hline
R-10 & 0.40 & — & +0.40 (미배정 풀에서) & 0.80 & 상승 \\ \hline
R-11 & 0.40 & — & −0.10 & 0.30 & 하락 \\ \hline
R-12 & 0.20 & — & −0.20 & 0 (이관) & 해제 \\ \hline
R-13 & 0.30 & — & 0 & 0.30 & 유지 \\ \hline
R-15 & — & — & +0.30 & 0.30 & 신규 \\ \hline
R-16 & — & — & +0.30 & 0.30 & 신규 \\ \hline
R-18 & — & — & +0.20 & 0.20 & 신규 \\ \hline
R-19 & — & — & +0.60 & 0.60 & 신규 (EX-01 대안 1 원가) \\ \hline
R-20 & — & — & +0.30 & 0.30 & 신규 \\ \hline
\textbf{배정 소계} & \textbf{4.60} & \textbf{1.07 / 1.07} & & \textbf{5.30} (13건, 잔여 EMV 컨틴전시 대상 5.06을 0.1 단위 반올림) & \\ \hline
미배정 풀 & 5.60 & \textbf{0.34 / 0.12} (CR-03 0.12 발생, CR-07 0.22 약정) & 승인 변경 BL-02 편입 −1.005 (CR-06 0.32 + CR-08 0.495 + CR-09 0.19), 배정 이동 −1.785 (신규 1.70 + R-10 증액 0.40 + R-05·R-07 재충전 0.47 − 해제 0.80 = 1.77, 반올림 +0.015) & \textbf{2.47} (스프린트 1.75회분) & Day2 정의: 스프린트 4회분 → 1.75회분 \\ \hline
\textbf{합계} & \textbf{10.20} & \textbf{1.41 / 1.19} & & 5.30 + 2.47 + 집행 1.41 + BL-02 편입 1.02 = \textbf{10.20} (검산 일치) & \\ \hline
관리예비비 요청 & \multicolumn{5}{>{\raggedright\arraybackslash}p{\dimexpr\linewidth-24mm-4\tabcolsep-3\arrayrulewidth\relax}|}{0건 (R-04 흡수, R-01 분리는 기준선 내). 잔여 12.0. 대상 리스크 EMV 2.20, 유보 복귀 조건 2건(클라우드 3.0 — IS-10 판단 2027-01, PMO 3.0 — G4 후)} \\ \hline
소진율 & \multicolumn{5}{>{\raggedright\arraybackslash}p{\dimexpr\linewidth-24mm-4\tabcolsep-3\arrayrulewidth\relax}|}{약정 1.41 / 10.20 = \textbf{13.8\%} vs 원가 진척 EV/BAC = 73.43 / 145.80 = \textbf{50.4\%}. 진척 대비 소진은 느림 — 그러나 배정 재조정 후 미배정 풀이 4회분 → 1.75회분: \textbf{다음 신규 리스크 2건이면 풀 소진}} \\ \hline
\fi
\end{longtable}}
\landscapeoff

% ------------------------------------------------------------------ 8 (세로)
\secbar{8. 예외 보고 필요 여부 판정과 에스컬레이션 현황}
\guide{리스크 축(잔여 EMV 합·단일)의 판정 한 줄 — 판정 없이 표만 있으면 안 된다. 관리예비비 대상 트리거 보고 이력. 이슈 에스컬레이션 단계별 건수와 응답 시한(3·10영업일) 준수. 이번 갱신에서 스폰서에 올리는 결정 요청은 ⑫ 성과보고의 결정 요청과 번호를 맞춘다.}
\begin{fieldbox}\small
\ifwbblank
\g{리스크 축 예외 보고: [불요/필요 — 잔여 EMV 합 [ ] vs 8, 단일 최대 [ ] vs 2, 경고 사유]}\lines{2}{7mm}\g{관리예비비 대상 트리거 보고 이력: [R-nn 일자(당일 여부)]}\lines{2}{7mm}
\else
\textbf{리스크 축 예외 보고: 불요.} 잔여 위협 EMV 7.26 $<$ 8.0, 단일 최대 0.80 $<$ 2.0. 경고 사유: 여유 0.74는 R-10(0.76) 하나의 크기이고, R-10의 트리거(정본 판정 후 변경 요청)는 G2 직후에 관측된다 — \textbf{11월 갱신에서 R-10이 실현되면 즉시 8억 초과}. 스폰서 사전 예고(10-05 성과보고 제9호).\par
관리예비비 대상 트리거 보고 이력: R-04 02-20(당일), R-01 04-24(당일, T0 통보 회의에서 구두 + 04-27 서면). R-14·R-17 신규 등재는 10-23 발행 시 보고.
\fi
\end{fieldbox}
\vspace{1mm}
{\footnotesize
\begin{tabularx}{\linewidth}{|L{52mm}|L{48mm}|L{30mm}|Y|}\hline
\hd{에스컬레이션 단계} & \hd{건수 (데이터 일자)} & \hd{응답 시한 준수} & \hd{미응답 · 지연} \\ \hline
\ifwbblank
\rd{8mm}1단계 (수행사 PM → P1 PM, 3영업일) & & & \\ \hline
\rd{8mm}2단계 (P1 PM → 스폰서, 10영업일) & & & \\ \hline
\rd{8mm}3단계 (스폰서 → 이사회) & & & \\ \hline
\rd{8mm}병행 (변경협의체 / 소관 임원 / 운영위) & & & \\ \hline
\else
1단계 (수행사 PM → P1 PM, 3영업일) & 3 (IS-03·10·11) & 평균 2영업일 — 준수 & — \\ \hline
2단계 (P1 PM → 스폰서, 10영업일) & 2 (IS-01 결정 03-06 준수, IS-07 회신 요청 09-30 → 10-08 대기) & 1/1 준수, 1 대기 & — \\ \hline
3단계 (스폰서 → 이사회) & 0 (R-01 분리 의결은 G1 정기 부의) & — & — \\ \hline
병행 (변경협의체 / 소관 임원 / 운영위) & 5 (IS-02·04·05·06·08) & 준수 & IS-05 CCB 보고 누락(8월) → 10-15 소급 \\ \hline
\fi
\end{tabularx}}
\vspace{1mm}\noindent\textbf{\small 결정 요청} \g{(⑫ 항목 10과 동일 번호)}\par\vspace{0.5mm}
{\footnotesize
\begin{tabularx}{\linewidth}{|C{6mm}|L{40mm}|L{30mm}|Y|}\hline
\hd{\#} & \hd{누가} & \hd{언제까지} & \hd{무엇을} \\ \hline
\ifwbblank
\rd{8mm}1 & & & \\ \hline
\rd{8mm}2 & & & \\ \hline
\rd{8mm}3 & & & \\ \hline
\rd{8mm}4 & & & \\ \hline
\else
1 & 정미란 · 프로그램 PMO장 & 2026-10-16 & EX-01 회복 계획 승인 · 합류 버퍼 배정 원칙(R-14 관련) 회신 \\ \hline
2 & CCB (정미란 위원장) & 2026-10-15 & CR-08~11 의결, 컨틴전시 8월 보고 누락 소급 확인 \\ \hline
3 & 정하늘 · 정미란 & 2026-10-08 & 감리 지적 3·4번 귀속 재분류 회신 \\ \hline
4 & 정미란 (재경팀장) & 2026-11 재무본부 초안 전 & R-17 — 2027 자금 계획에 이월분 12.7(4.70 + 8.0) 반영 시점 확인 \\ \hline
\fi
\end{tabularx}}

% ------------------------------------------------------------------ 9
\secbar[50mm]{9. 검토 주기와 다음 갱신}
\guide{다음 갱신 시점(월·게이트), 게이트에서 정식 재평가할 항목, 소유자 서명 현황, 정량 분석(몬테카를로) 실시 여부. 다음 갱신을 정하지 않은 등록부는 갱신되지 않는다.}
\begin{fieldbox}\small
\ifwbblank
\g{다음 월 갱신: [일자 — 데이터 일자·회의]}\lines{1}{7mm}\g{게이트 재평가: [G-n 일자 — 재평가 항목 R-nn과 판단 근거]}\lines{2}{7mm}\g{소유자 서명: [n건, 시점 — 미서명 사유]}\lines{1}{7mm}\g{정량 분석: [몬테카를로 실시 여부·주체 / 시나리오 분석 제출]}\lines{2}{7mm}
\else
• 월 갱신: 2026-10-31 데이터(FZ-10 10-28 다음 주 PM 회의 11-04), 신규 접수 요구와 함께 신규 리스크 등록.\par
• G2 재평가(2026-11-27): R-10(정본 판정·동결 규칙 의결 결과), R-14(2차 컷오버 계획 확정, P2 MC 재통보), R-02(배포 범위 확정), R-19(S22 속도 판정 11-13 결과), R-17(2027 자금 계획 초안).\par
• 소유자 서명: v1.2 발행 시 22건(위협 16 + 기회 1 + 종결·이관 5의 종결 확인) — R-14 소유자 프로그램 PMO장 서명은 프로그램 운영위 11월.\par
• 정량 분석: 3점 추정 몬테카를로는 G2에서 프로그램 PMO 판단. P1은 R-14 × R-19 결합 시나리오(두 번 컷오버 중 결함 유출)의 시나리오 분석표(3안)를 G2 자료로 제출 [추가 설정].
\fi
\end{fieldbox}
\end{document}
